% =============================================================
% §2 背景
% =============================================================
\section{背景}
\label{sec:background}

\subsection{LLM 推理：Prefill、Decode 与 KV cache}
\label{sec:bg-inference}

\begin{figure*}[t]
  \centering
  \includegraphics[width=0.70\textwidth]{llm_inference_lifecycle.pdf}
  \caption{一次 LLM 推理请求的两阶段生命周期。\emph{Prefill} 对整段 prompt 执行一次（计算受限），输出首个 token 并建立 KV cache；\emph{Decode} 随后执行 $L_o$ 个串行的、显存带宽受限的步骤，每步读取并追加 KV cache。KV cache 绑定在运行 Prefill 的那台 GPU 上，因此把 Decode 派到别处会强制 KV 迁移或重跑 Prefill——这正是跨服务器调度器必须权衡的核心本地性张力。}
  \label{fig:lifecycle}
\end{figure*}

现代 LLM 是 decoder-only Transformer~\cite{vaswani2017attention,brown2020gpt3}，其推理分为两个硬件特征截然不同的阶段（图~\ref{fig:lifecycle}）。\textbf{Prefill} 每请求执行一次：在一次使算力饱和的前向传播中处理长度为 $L_p$ 的 prompt，产生首个输出 token 及每一层对应的\emph{键值缓存}（KV cache）。\textbf{Decode} 随后迭代执行，每个输出 token 一步（A100 上约 $10$--$80$~ms），复用并追加 KV cache；对长生成场景，$L_o$ 步序列主导端到端延迟。

KV cache 体积大（每 token 每层 $0.5$--$2.5$~MB；1000-token 的 LLaMA-13B 会话约 800~MB）且\textbf{与运行 Prefill 的 GPU 绑定}：在别处服务 Decode 要么通过网络拷贝数百 MB，要么重跑 Prefill。

除 Transformer LLM 之外，AIGC 负载还包括文本生成图像的扩散模型，其推理在每个去噪步都调用同一个共享 U-Net——在结构上与 LLM 不同，但共享同样的模型驻留与批处理特性，这正是本文工作的动机所在。

\subsection{连续批处理与现代服务系统}
\label{sec:bg-batching}

\textbf{连续（迭代级）批处理}——由 Orca~\cite{yu2022orca} 引入、vLLM~\cite{kwon2023vllm}（PagedAttention）与 Sarathi-Serve~\cite{agrawal2024sarathi}（分块 prefill）进一步精炼——在每个 decoder 迭代都重组当前批：完成的请求即时退出，排队请求中途加入。每请求批处理开销在 decode 阶段 $<5\%$（显存带宽受限），在 prefill 阶段约 30\%（算力受限）——即本文仿真器采用的阶段相关系数（\S\ref{sec:simulator}）。这些引擎级优化仅作用于\textbf{单台服务器之内}；\emph{某个请求由哪台服务器处理}这一上游问题正是本文聚焦的对象。

\subsection{云--边协同 LLM 部署}
\label{sec:bg-cloudedge}

生产环境的 LLM 服务越来越多地跨越异构算力池部署：一台或多台云服务器配备高端加速器（A100/H100，GPU 显存即 VRAM $\ge$80~GB），可承载 LLaMA-70B-INT8（约 70~GB）的高吞吐服务；一组边缘服务器涵盖数据中心 A100、工作站级 T4 直至嵌入式 Jetson Orin，显存 16--64~GB、算力 10--50 TFLOPS、热设计功耗（TDP）30--250~W，适合 LLaMA-7B/13B 等较小变体；以及一张云--边往返延迟 30--65~ms、边--边往返延迟 8--23~ms 的网络。该范式由\textbf{延迟}（边缘临近降低首 token 延迟，TTFT）、\textbf{成本}（边缘 GPU 每 FLOP 便宜 1--10$\times$，功耗低 2--10$\times$）与\textbf{数据主权}驱动。阶段解耦与迁移类系统（\S\ref{sec:related}）重新设计部署本身；我们回答它们留下的互补的\emph{初始放置}问题：给定固定的异构池，每个请求应派往何处？

\subsection{调度方法}
\label{sec:bg-sched}

云--边调度历来使用启发式调度器（HEFT~\cite{topcuoglu2002heft}、Round-Robin/Least-Loaded/Shortest-Queue）、元启发式（GA~\cite{holland1975adaptation}、PSO~\cite{kennedy1995pso}）与学习式调度器（Decima~\cite{mao2019decima}、RLScheduler~\cite{zhang2020rlscheduler}，以及拓扑上最接近的 A3C-R2N2~\cite{tuli2022a3cr2n}——面向通用边缘--云任务、采用门控循环单元（GRU）+ 残差编码器的异步优势 actor-critic）。它们均未原生编码 AIGC 物理：启发式将任务视作独立黑盒，元启发式的 fitness 函数需要手工设计 AIGC 项——这正是本文自动化的环节——而现有学习式调度器均早于 LLM 时代。我们在第~\ref{sec:related} 节深入综述这一图景，并在第~\ref{sec:eval} 节对三类方法全部进行基准比较。
